
Este capítulo presenta los resultados de los cinco experimentos descritos en
el Capítulo \ref{CAP:diseno} y los discute en su conjunto. Los resultados toman como dominio
principal CIFAR-10 y se contrastan, cuando añaden información, con SVHN, CIFAR-100
y FashionMNIST.

% ═══════════════════════════════════════════════════════════════════════════════
\section{Resultados de Accuracy: Búsqueda de Hiperparámetros (EXP-1)\label{sec:resultados_accuracy}}

Para cada técnica se entrena un conjunto de modelos mediante \textit{grid search} sobre
el hiperparámetro de regularización; para cada método se retiene el modelo con mayor
\textit{val\_acc} al final de las 60 épocas de entrenamiento sobre CIFAR-10.

Los resultados revelan tres grupos con comportamientos claramente diferenciados
(Figura~\ref{fig:ranking_cifar10}). \textbf{Data Augmentation} y \textbf{Batch
Normalization} logran mejoras sustanciales sobre el baseline, siendo los únicos métodos
que superan 0.75 de val\_acc. En un nivel intermedio se encuentran \textbf{Dropout}
y \textbf{L2}, con mejoras moderadas. El tercer grupo (\textbf{Gaussian Noise},
\textbf{Early Stopping} y especialmente \textbf{L1}) apenas supera o ni siquiera
alcanza el baseline: L1 es el único método que perjudica el rendimiento, reflejo de que
su ventana de hiperparámetro viable es extremadamente estrecha y el modelo colapsa
abruptamente para $\lambda \gtrsim 10^{-3}$.

\begin{figure}{fig:ranking_cifar10}{Ranking de val\_acc final en CIFAR-10 (izquierda) y
    SVHN (derecha). El orden relativo entre métodos se preserva en ambos dominios;
    el rendimiento absoluto es mayor en SVHN por la mayor homogeneidad de sus imágenes.}
    \begin{minipage}{0.49\textwidth}
        \includegraphics[width=\textwidth]{fig02_ranking_val_acc_CIFAR10}
    \end{minipage}
    \hfill
    \begin{minipage}{0.49\textwidth}
        \includegraphics[width=\textwidth]{fig02_ranking_val_acc_SVHN}
    \end{minipage}
\end{figure}

Más allá del ranking, la \textbf{brecha de generalización} (\textit{gap} = train\_acc
$-$ val\_acc) diferencia nítidamente a los métodos. Baseline, L2 y Gaussian Noise
alcanzan train\_acc~$\approx 1.0$ con un gap próximo a 0.30: el modelo memoriza el
conjunto de entrenamiento sin que ello se traduzca en mejora de val\_acc. DataAug, en
cambio, reduce la brecha a~$\sim$0.09 con una train\_acc de~0.877, evidenciando
regularización implícita en el espacio de datos. BatchNorm logra alta val\_acc pero
mantiene un gap de~$\sim$0.23, lo que sugiere que su ventaja proviene principalmente de
estabilizar la optimización, no de suprimir el sobreajuste. Un caso particular es Early
Stopping: detiene el entrenamiento en torno a la época~11 (patience=8) y aun así alcanza
val\_acc competitiva, lo que apunta a que una fracción importante del sobreajuste en
esta arquitectura se concentra en las épocas tardías del entrenamiento.

\begin{figure}{fig:curvas_cifar10}{Análisis de Data Augmentation en CIFAR-10 (EXP-1).
    Arriba: mapas de calor de val\_acc y train\_acc a lo largo de las 60 épocas para
    cada nivel de aumentación. Abajo: mejor val\_acc y brecha de generalización en función
    del nivel. La brecha desciende de $\sim$0.31 a $\sim$0.05, mostrando el efecto
    regularizador más claro del experimento.}
    \includegraphics[width=0.90\textwidth]{fig01_accuracy_curvas_CIFAR10}
\end{figure}

El panel derecho de la Figura~\ref{fig:ranking_cifar10} muestra el mismo experimento
sobre SVHN. El rendimiento absoluto aumenta (las imágenes de dígitos son más homogéneas
que las de CIFAR-10) pero el \textbf{ranking relativo se preserva}: el orden entre
métodos es prácticamente idéntico al de CIFAR-10. En CIFAR-100, la complejidad del
problema supera la capacidad de esta arquitectura y los rendimientos caen de forma
drástica para todos los métodos: DataAug sigue liderando (val\_acc\,=\,0.44) seguido
de BatchNorm (0.42), pero el resto queda agrupado entre 0.33 y 0.36, y el cuello de
botella es la capacidad del modelo, no el tipo de regularización. En FashionMNIST se
observa el efecto contrario: todos los métodos superan 0.91 de val\_acc y las diferencias
entre ellos no alcanzan el 0.01, lo que constituye un techo de rendimiento claro para
esta arquitectura sobre esta tarea. Llama la atención que el ranking invierte
parcialmente: BatchNorm y Dropout lideran, por delante de DataAug, a diferencia de lo
observado en CIFAR-10 y CIFAR-100. Este patrón transversal (conservación del ranking en
SVHN, compresión hacia el techo en FashionMNIST, colapso de capacidad en CIFAR-100)
confirma que CIFAR-10 constituye el punto de operación más informativo del experimento
y que las conclusiones no son artefactos específicos del dominio.

\FloatBarrier
% ═══════════════════════════════════════════════════════════════════════════════
\section{Análisis de Activaciones Internas\label{CAP:resultados_activaciones}}

Los resultados de \textit{accuracy} miden el efecto externo de la regularización.
A continuación se adopta la perspectiva interna: cómo cada técnica altera la
distribución y diversidad de las representaciones ocultas de la red.

\subsection{Correlación entre Métricas Internas y Rendimiento (EXP-2)\label{sec:correlacion_metricas}}

El análisis de activaciones caracteriza la estructura interna de las representaciones
a través de tres métricas complementarias. El \textbf{porcentaje de estados únicos
(\textit{unique\_pctg})} mide la fracción de ejemplos de validación que producen una
activación discreta diferente tras cuantizar el histograma en \textit{bins}: captura
directamente la capacidad del modelo de distinguir representaciones individuales. La
\textbf{entropía de Shannon} ($H$) del histograma cuantifica la uniformidad de la
distribución de activaciones. El \textbf{ratio de dispersión} cuantifica la amplitud
de las activaciones relativa al baseline. Aunque las tres miden aspectos distintos de la
compacidad representacional, \textit{unique\_pctg} es especialmente útil para distinguir
entre métodos, ya que captura directamente si el modelo colapsa ejemplos distintos a un
mismo estado cuantizado.

El estudio se centra en \texttt{fc1}, la última capa oculta antes del clasificador.
Teóricamente, es la capa que agrega las características convolucionales en
representaciones semánticas de clase, y su estructura determina directamente la
discriminabilidad final del modelo. Empíricamente, calculando la desviación típica de
\textit{unique\_pctg} sobre todos los valores del hiperparámetro de cada método,
\texttt{fc1} muestra una variabilidad media de 27.2~pp frente a 10.9~pp de \texttt{conv3}
y 2.6~pp de \texttt{conv1}: las capas tempranas aprenden filtros locales estables,
independientes de la regularización, mientras que \texttt{fc1} es la capa diagnóstica
por excelencia.

Cabe señalar que los experimentos emplean tanto 10 como 30 \textit{bins} para CIFAR-10 y para SVHN. Los
valores absolutos de \textit{unique\_pctg} no son directamente comparables entre
resoluciones, pero la \textbf{estructura de ordenamiento relativo} entre métodos se
preserva. La Tabla~\ref{tab:metricas_internas} lo ilustra mostrando ambas resoluciones
en el escenario óptimo.

\begin{table}{tab:metricas_internas}{Métricas de \textit{unique\_pctg} y entropía $H$ en \texttt{fc1}
para la configuración óptima de cada método (CIFAR-10, validación).
La comparativa 30 vs.\ 10 \textit{bins} muestra que el ordenamiento relativo entre
métodos es invariante a la resolución de discretización.}
\small
\begin{tabular}{lc|cc|cc}
\toprule
& & \multicolumn{2}{c|}{\textbf{30 bins}} & \multicolumn{2}{c}{\textbf{10 bins}} \\
\textbf{Método} & \textbf{val\_acc} & \textbf{uniq\_pctg (\%)} & \textbf{$H$}
  & \textbf{uniq\_pctg (\%)} & \textbf{$H$} \\
\midrule
Data Aug    & 0.785 &  89.22 & 1.78 &  7.12 & 0.49 \\
Batch Norm  & 0.770 &  84.30 & 1.53 &  4.29 & 0.39 \\
Dropout     & 0.727 &  37.23 & 1.17 &  2.83 & 0.08 \\
L2          & 0.714 &  97.93 & 2.09 & 11.70 & 0.72 \\
Gauss Noise & 0.697 & 100.00 & 2.93 & 68.35 & 1.52 \\
Early Stop  & 0.696 &  20.56 & 0.88 &  0.63 & 0.03 \\
Baseline    & 0.694 &  99.98 & 2.87 & 61.41 & 1.47 \\
L1          & 0.694 &  99.98 & 2.87 & 61.41 & 1.47 \\
\bottomrule
\end{tabular}
\end{table}

La Tabla~\ref{tab:metricas_3dim} presenta las tres métricas para el
escenario óptimo de cada método, confirmando la tesis del \textit{sweet spot} desde tres
puntos de vista: los tres métodos con mejor val\_acc (DataAug, BatchNorm, Dropout) presentan
\textit{unique\_pctg} intermedio, $H$ intermedia \textit{y} ratio de dispersión menor
que la unidad (representaciones más compactas que el baseline). Los métodos con peor
rendimiento incumplen al menos uno de los tres criterios: \textit{unique\_pctg} o $H$
en la zona saturada ($\approx$100\% / $\approx$2.9), o ratio de dispersión $\geq 1.0$.

\begin{table}{tab:metricas_3dim}{Las tres métricas de activación en \texttt{fc1} para
el escenario óptimo de cada método (CIFAR-10, 30 bins). Los métodos en la zona de
alto rendimiento satisfacen simultáneamente los tres criterios de la zona intermedia.}
\small
\begin{tabular}{lcccc}
\toprule
\textbf{Método} & \textbf{val\_acc} & \textbf{uniq\_pctg (\%)} & \textbf{Entropía $H$}
  & \textbf{Disp.\ ratio} \\
\midrule
Data Aug    & 0.785 &  89.22 & 1.78 & 0.46 \\
Batch Norm  & 0.770 &  84.30 & 1.53 & 0.41 \\
Dropout     & 0.727 &  37.23 & 1.17 & 0.29 \\
L2          & 0.714 &  97.93 & 2.09 & 0.57 \\
Gauss Noise & 0.697 & 100.00 & 2.93 & 1.11 \\
Early Stop  & 0.696 &  20.56 & 0.88 & 0.22 \\
Baseline    & 0.694 &  99.98 & 2.87 & $\sim$1.00 \\
L1          & 0.694 &  99.98 & 2.87 & 1.06 \\
\bottomrule
\end{tabular}
\end{table}

Un aspecto clave de la Tabla~\ref{tab:metricas_3dim} es la \textbf{covarianza positiva
entre las tres métricas}: los métodos en la zona saturada (GaussNoise, Baseline, L1)
presentan simultáneamente \textit{unique\_pctg} alto, $H$ alta y ratio de dispersión
$\geq$1.0; los métodos en la zona de colapso (EarlyStopping, Dropout sobre-regularizado)
muestran las tres métricas bajas al mismo tiempo. Esta co-varianza no es accidental: las
tres son facetas distintas de la misma cantidad subyacente (la diversidad informacional
del espacio de activación). Que el \textit{sweet spot} aparezca de forma coherente en
las tres medidas simultáneamente refuerza la robustez empírica del hallazgo.

Los datos revelan además un \textbf{patrón no lineal} confirmado sobre el \textit{grid
search} completo: las configuraciones con \textit{unique\_pctg} intermedio
($\approx$30--90\% a 30 bins) concentran los mejores val\_acc, mientras los extremos
(saturación de $\approx$100\% o colapso por debajo del 25\%) corresponden a peor
generalización. Las Figuras~\ref{fig:entropia_vs_acc} y~\ref{fig:dispersion_vs_acc}
visualizan este patrón sobre \textbf{todas} las configuraciones evaluadas (no solo las
óptimas), con el eje de val\_acc restringido a la región de interés para resaltar la
curvatura. L2 es una excepción parcial: su óptimo cae en unique\_pctg~$\approx$~98\%
pero el ratio de dispersión permanece por debajo de la unidad (0.57), porque la
restricción $\ell_2$ actúa sobre la magnitud de los pesos sin necesariamente comprimir
las representaciones intermedias.

\begin{figure}{fig:entropia_vs_acc}{Dispersión de val\_acc frente a \textit{unique\_pctg}
    en \texttt{fc1} (CIFAR-10, 30 bins, todas las configuraciones del \textit{grid search}).
    La zona intermedia de \textit{unique\_pctg} concentra los mejores val\_acc.}
    \includegraphics[width=0.82\textwidth]{fig03_entropia_vs_acc}
\end{figure}

\begin{figure}{fig:dispersion_vs_acc}{Dispersión de val\_acc frente al ratio de
    dispersión en \texttt{fc1} (CIFAR-10). Los métodos mejor regularizados muestran
    dispersión inferior a la del baseline (ratio~$<$~1).}
    \includegraphics[width=0.82\textwidth]{fig04_dispersion_vs_acc}
\end{figure}

Una comparativa especialmente ilustrativa es la de L1 frente a L2: con val\_acc
similares, L2 mantiene representaciones ricas (\textit{unique\_pctg}~$\approx$~98\%),
mientras que L1 en su rango regularizado las colapsa drásticamente. Dos regularizadores
con rendimiento comparable actúan por mecanismos internos radicalmente distintos.

El mismo estudio replicado en SVHN (Figura~\ref{fig:svhn_unique_pctg_vs_acc}) muestra que el
patrón tripartito (sub-regularización, zona óptima, colapso) se reproduce con coherencia
al de CIFAR-10, confirmando la validez diagnóstica de la métrica entre dominios.

En CIFAR-100, sin embargo, la señal interna desaparece por saturación: todos los
métodos convergen a \textit{unique\_pctg} cercano al 100\% en \texttt{fc1}, con entropías
en el rango $[2.1,\,3.5]$, muy por encima del intervalo óptimo de CIFAR-10. La causa
es la mayor dimensionalidad del problema (100 clases): la red necesita representaciones
más diversas para separar todas las categorías, y la zona de colapso simplemente no
aparece en el rango de hiperparámetros evaluado. En este contexto, la métrica interna
deja de ser discriminante porque el cuello de botella es la capacidad de la arquitectura,
no la intensidad de la regularización.

En FashionMNIST (Figura~\ref{fig:fashionmnist_unique_pctg_vs_acc}) el comportamiento
es distinto. La mayoría de métodos satura igualmente cerca del 100\%, pero DataAug
se sitúa en el 73.6\% (dentro de la zona óptima identificada en CIFAR-10) mientras
mantiene un \textit{val\_acc} competitivo. El efecto techo de rendimiento (todos los
métodos entre 0.91 y 0.92) diluye la correlación con la métrica, pero la posición de
DataAug ilustra que, incluso en tareas sencillas, una regularización basada en diversidad
de datos produce representaciones con menor saturación en \texttt{fc1}.

\begin{figure}{fig:svhn_unique_pctg_vs_acc}{Dispersión de val\_acc frente a
    \textit{unique\_pctg} en \texttt{fc1} para SVHN (10 bins). El patrón tripartito
    es análogo al de CIFAR-10, confirmando la consistencia de la métrica entre dominios.}
    \includegraphics[width=0.82\textwidth]{fig10_svhn_unique_pctg_vs_acc}
\end{figure}

\begin{figure}{fig:fashionmnist_unique_pctg_vs_acc}{Dispersión de val\_acc frente a
    \textit{unique\_pctg} en \texttt{fc1} para FashionMNIST (30 bins). La mayoría de
    métodos se agrupa cerca del 100\% por efecto techo; DataAug se mantiene en la zona
    moderada (73.6\%), coherente con el rango óptimo identificado en CIFAR-10.}
    \includegraphics[width=0.82\textwidth]{fig11_fashionmnist_unique_pctg_vs_acc}
\end{figure}

\FloatBarrier
% ─────────────────────────────────────────────────────────────────────────────
\subsection{Evolución Temporal de las Representaciones (EXP-3)}

La Figura~\ref{fig:temporal_custom} muestra la evolución de \textit{unique\_pctg} durante
las 60 épocas de entrenamiento en las tres capas: \texttt{conv1}, \texttt{conv3} y
\texttt{fc1}.

\begin{figure}{fig:temporal_custom}{Evolución temporal de \textit{unique\_pctg} en las
    tres capas de la CNN (CIFAR-10, serie Custom). \texttt{conv1} permanece al 100\%
    en todos los métodos durante todo el entrenamiento. \texttt{conv3} y \texttt{fc1}
    muestran dinámicas muy distintas: la compresión inicial de \texttt{fc1} y su
    recuperación diferenciada entre métodos es el principal elemento diagnóstico.}
    \includegraphics[width=0.96\textwidth]{fig05_temporal_uniquePctg_Custom}
\end{figure}

Los datos de la época~1 revelan un contraste muy marcado entre capas. \texttt{conv1}
alcanza el 100\% de \textit{unique\_pctg} para \textbf{todos} los métodos desde el
inicio: esta capa opera directamente sobre los parches de píxeles de las imágenes de
entrada, que son intrínsecamente diversos entre sí; la alta dimensionalidad de entrada
y el receptivo campo local garantizan representaciones únicas por imagen incluso con
pesos aleatorios. En el extremo opuesto, \texttt{fc1} comienza prácticamente colapsada:
Baseline, GaussNoise y L2 arrancan con unique\_pctg~$\approx$~1\%, Dropout y DataAug
con~$\approx$~0.5\%. El motivo es geométrico: \texttt{fc1} es una proyección lineal
totalmente conectada que agrega toda la información espacial post-pooling en un único
vector; con pesos aleatorios, esta proyección de alta a baja dimensionalidad colapsa
entradas distintas a activaciones discretas muy similares. \texttt{conv3} arranca en
una posición intermedia (4--12\%), ya parcialmente comprimida por los pasos de
\textit{pooling} previos. La excepción es \textbf{BatchNorm}, que en la primera época
muestra \texttt{conv3}~$\approx$~67\% y \texttt{fc1}~$\approx$~42\%: la normalización
por lotes regulariza las activaciones desde el primer paso hacia delante e impide el
colapso inicial.

El seguimiento a lo largo de las 60 épocas permite identificar tres fases con dinámicas
distintas en \texttt{fc1}. Durante las primeras diez épocas \texttt{fc1} permanece muy
comprimida (2--15\%) mientras \texttt{conv3} se recupera con rapidez (67--98\%): la red
aprende filtros básicos y la información todavía no fluye de forma discriminativa hacia
las capas superiores, mientras \texttt{conv1} permanece inalterada, ya saturada desde el
inicio. Entre las épocas 10 y 30 se produce la divergencia crítica entre métodos:
\textbf{Baseline} y \textbf{Gaussian Noise} recuperan \texttt{fc1} de forma acelerada,
alcanzando $\approx$99\% al final del periodo, mientras que \textbf{DataAug} y
\textbf{Dropout} mantienen la compresión en torno al 40\% (la aumentación de datos y
el apagado aleatorio de neuronas obligan a construir representaciones más compactas) y
\textbf{L2} también permanece comprimido ($\approx$41\%) por la penalización continua
sobre los pesos. A partir de la época 30 el sistema se estabiliza: \textbf{Baseline}
y \textbf{Gaussian Noise} saturan al 100\%,
sin compresión representacional residual. \textbf{L2} termina en el 44.8\%, el valor
más bajo de la serie, señal de una compresión continua y activa. \textbf{Batch Norm}
y \textbf{Dropout} convergen a valores intermedios (85--88\%), coherentes con la zona
de buen rendimiento identificada en EXP-2. \textbf{DataAug} termina en el 95.7\% en
la serie Custom, pero en el escenario óptimo (nivel de aumentación 2) su \texttt{fc1}
alcanza 89.2\%, confirmando que el nivel de aumentación controla directamente el grado
de compresión representacional en la capa superior.

\FloatBarrier
% ─────────────────────────────────────────────────────────────────────────────
\subsection{Robustez ante Ruido en los Datos (EXP-4)}

Se añade ruido gaussiano de intensidad creciente a las imágenes de validación,
manteniendo los pesos inalterados, para evaluar cuánto rendimiento retiene cada método.

\begin{figure}{fig:robustez_datos}{Curvas de degradación del \textit{accuracy} ante ruido gaussiano en los
    datos de entrada (CIFAR-10, serie Custom). Los modelos fueron entrenados sin ruido.}
    \includegraphics[width=0.82\textwidth]{fig06_robustez_datos_Custom}
\end{figure}

\begin{table}{tab:robustez_datos}{Retención de \textit{accuracy} ante ruido máximo en datos de entrada
(CIFAR-10, serie Custom, $N_{\text{RUNS}}$=3).}
\begin{tabular}{lcc}
\toprule
\textbf{Método} & \textbf{acc($\sigma$=0)} & \textbf{Retención a $\sigma$=0.5 (\%)} \\
\midrule
Data Aug        & 77.7\% & 18.9 \\
Batch Norm      & 76.8\% & 21.1 \\
Dropout         & 71.2\% & 23.9 \\
L2              & 68.2\% & 31.2 \\
Baseline        & 68.5\% & 33.3 \\
\textbf{Gauss Noise} & 63.9\% & \textbf{51.6} \\
\bottomrule
\end{tabular}
\end{table}



Los resultados revelan una \textbf{disyuntiva entre rendimiento y robustez a entradas
ruidosas}: los métodos con mayor \textit{accuracy} en condiciones limpias son
precisamente los que peor retienen su rendimiento ante ruido intenso. Gaussian Noise
es el método más robusto ante este tipo de perturbación (directamente como consecuencia
de haber entrenado con ruido del mismo tipo), aunque parte de un \textit{accuracy}
base menor. La regularización orienta el aprendizaje hacia invarianzas específicas al
tipo de variabilidad del entrenamiento, no hacia una robustez genérica.

\FloatBarrier
% ─────────────────────────────────────────────────────────────────────────────
\subsection{Robustez ante Ruido en los Pesos y Análisis de Mínimos Planos (EXP-5)}

Se añade ruido gaussiano a los pesos de capas específicas tras el entrenamiento para
medir la estabilidad paramétrica de cada modelo.

\begin{figure}{fig:robustez_conv_layers}{Degradación del \textit{accuracy} al perturbar los pesos de
    \texttt{conv1} (izquierda) y \texttt{conv3} (derecha) (CIFAR-10, Custom).}
    \begin{minipage}{0.49\textwidth}
        \includegraphics[width=\textwidth]{fig07_robustez_pesos_conv1_Custom}
    \end{minipage}
    \hfill
    \begin{minipage}{0.49\textwidth}
        \includegraphics[width=\textwidth]{fig07b_robustez_pesos_conv3_Custom}
    \end{minipage}
\end{figure}

El hallazgo más llamativo es \textbf{BatchNorm}: a pesar de ser el segundo método más
preciso, colapsa drásticamente ante perturbaciones en los pesos de las capas
convolucionales, indicando que converge a mínimos <<agudos>> (\textit{sharp minima}).

\begin{figure}{fig:robustez_fc1}{Degradación del \textit{accuracy} al perturbar los pesos de \texttt{fc1}
    (CIFAR-10, Custom). Dropout muestra la mayor robustez, coherente con la búsqueda de
    mínimos planos.}
    \includegraphics[width=0.75\textwidth]{fig08_robustez_pesos_fc1_Custom}
\end{figure}

\begin{table}{tab:robustez_pesos}{Retención de \textit{accuracy} a $\sigma$=0.15 de ruido en pesos, por
capa (CIFAR-10, serie Custom).}
\begin{tabular}{lccc}
\toprule
\textbf{Método} & \textbf{Ret. Conv1 (\%)} & \textbf{Ret. Conv3 (\%)}
  & \textbf{Ret. FC1 (\%)} \\
\midrule
Gaussian Noise (0.1)   & \textbf{94.9} & \textbf{80.3} & 68.5 \\
Baseline               & 93.7 & 75.7 & 68.4 \\
Dropout (0.3)          & 91.0 & 79.4 & \textbf{79.9} \\
L2 (0.001)             & 89.5 & \textbf{40.7} & \textbf{39.9} \\
Data Aug (2.0)         & 83.8 & 67.5 & 69.9 \\
\textbf{Batch Norm}    & \textbf{41.0} & 41.4 & 60.6 \\
\bottomrule
\end{tabular}
\end{table}


\textbf{Dropout} es el único método con robustez alta y \textit{consistente} en las
tres capas. \textbf{L2} muestra un perfil heterogéneo: muy robusto en \texttt{conv1}
pero con colapso pronunciado en \texttt{conv3} y \texttt{fc1}, revelando una protección
que no se extiende uniformemente a las capas profundas. La
Figura~\ref{fig:comparativa_capas} resume visualmente esta heterogeneidad inter-capa.

\begin{figure}{fig:comparativa_capas}{Retención de \textit{accuracy} a $\sigma=0.15$ para cada método y capa.
    La consistencia de Dropout en todas las capas contrasta con el comportamiento
    heterogéneo de L2 y el colapso de BatchNorm en capas convolucionales.}
    \includegraphics[width=0.82\textwidth]{fig07c_comparativa_capas_sigma15_Custom}
\end{figure}

\FloatBarrier
% ─────────────────────────────────────────────────────────────────────────────
\subsection{Síntesis de los Experimentos}

El análisis conjunto revela una \textbf{tensión fundamental} entre accuracy, robustez
distribucional y estabilidad paramétrica. Ningún método domina simultáneamente en las
tres dimensiones:

\begin{table}{tab:resumen_dimensiones}{Mejor y peor método en cada dimensión de evaluación (CIFAR-10).}
\small
\begin{tabular}{lll}
\toprule
\textbf{Dimensión} & \textbf{Mejor} & \textbf{Peor} \\
\midrule
Accuracy                     & DataAug         & L1 (por debajo del baseline) \\
Robustez datos               & Gaussian Noise  & Batch Norm \\
Robustez pesos Conv1         & Gaussian Noise  & Batch Norm \\
Robustez pesos Conv3         & Gaussian Noise  & L2 \\
Robustez pesos FC1           & Dropout         & L2 \\
Consistencia inter-capas     & Dropout         & L2 \\
\bottomrule
\end{tabular}
\end{table}

\FloatBarrier
% ═══════════════════════════════════════════════════════════════════════════════
\section{Discusión\label{CAP:discusion}}

Los cinco experimentos anteriores aportan resultados desde perspectivas independientes:
\textit{accuracy} comparado, representaciones internas, evolución temporal, robustez
a ruido en datos y estabilidad paramétrica. A continuación se conectan estos resultados
entre sí y con la literatura existente.

El resultado más transversal del estudio es que \textbf{ningún método domina en todas
las dimensiones relevantes}. Si el \textit{accuracy} en validación fuera el único
criterio, la respuesta sería trivial: usar Data Augmentation. Los experimentos de
robustez revelan sin embargo una imagen considerablemente más compleja. Data Augmentation,
el método más preciso, es uno de los menos robustos ante ruido de píxel; Batch
Normalization converge a mínimos agudos en capas convolucionales y colapsa ante
perturbaciones moderadas de pesos; Gaussian Noise es el más robusto ante entradas
ruidosas pero sacrifica rendimiento; y Dropout es el único con robustez consistente
en todas las capas, convergiendo sistemáticamente a mínimos planos en línea con
resultados previos \cite{hochreiter1997flat,keskar2017large,wan2013regularization,gal2016dropout}.
L2, por su parte, presenta un perfil asimétrico: muy robusto en capas tempranas pero
frágil en capas profundas. La Tabla~\ref{tab:evaluacion_multidim} y la
Figura~\ref{fig:radar} sintetizan esta heterogeneidad, cuya implicación práctica directa
es que la elección del regularizador debería guiarse por las prioridades del problema,
no por un único indicador de rendimiento.

\begin{table}{tab:evaluacion_multidim}{Evaluación comparativa de los regularizadores en tres dimensiones:
\textit{accuracy}, robustez a ruido en datos y estabilidad paramétrica (CIFAR-10).}
\small
\begin{tabular}{lcccc}
\toprule
\textbf{Método} & \textbf{val\_acc} & \textbf{Rob.~datos} & \textbf{Rob.~pesos (fc1)} & \textbf{Mínimo} \\
\midrule
Gauss Noise & 0.697 & muy alta  & alta       & Plano-medio \\
L2          & 0.714 & media     & baja       & Heterogéneo \\
Baseline    & 0.694 & media     & alta       & Variable    \\
Dropout     & 0.727 & baja      & muy alta   & \textbf{Plano} \\
Data Aug    & 0.785 & muy baja  & alta       & Plano-medio \\
Batch Norm  & 0.770 & muy baja  & media      & \textbf{Agudo (conv)} \\
\bottomrule
\end{tabular}
\end{table}

\begin{figure}{fig:radar}{Evaluación multidimensional normalizada de los regularizadores (CIFAR-10). Ningún método
    domina en las tres dimensiones simultáneamente.}
    \includegraphics[width=0.68\textwidth]{fig09_radar_multidimensional}
\end{figure}

Esta heterogeneidad entre métodos no es, sin embargo, arbitraria. Los experimentos de
activaciones internas muestran que \textit{unique\_pctg}, entropía de Shannon y ratio
de dispersión en \texttt{fc1} muestran correlaciones sistemáticas con el rendimiento y
covariación positiva entre ellas. Existe una zona óptima de complejidad representacional
(entropía intermedia, diversidad moderada, ratio inferior a la unidad) donde se
concentran los métodos de mayor \textit{accuracy}; los métodos colapsados (Early Stopping)
o saturados (Gaussian Noise, Baseline) quedan desplazados hacia los extremos en las tres
métricas simultáneamente. Esta zona es interpretable: por debajo de ella la red pierde
poder discriminativo; por encima, tiende a memorizar patrones de entrenamiento. Desde el
punto de vista práctico, estas señales son detectables sin necesidad de esperar al
\textit{accuracy} final, lo que abre la puerta a diagnósticos tempranos durante el
desarrollo del modelo.

El experimento de evolución temporal añade una tercera dimensión al cuadro. Todos los
métodos, salvo BatchNorm, inician con una compresión representacional extrema en
\texttt{fc1} y se diferencian progresivamente a lo largo del entrenamiento, en un patrón
coherente con la hipótesis del \textit{information bottleneck}
\cite{tishby2017opening}, aunque medido aquí mediante una aproximación discreta.
BatchNorm evita este colapso inicial gracias al suavizado del paisaje de pérdida
\cite{santurkar2018does} y diferencia representaciones antes, pero a costa de mínimos más
agudos: los hallazgos de robustez paramétrica y de dinámica temporal resultan así
coherentes entre sí y se refuerzan mutuamente. La capacidad del baseline de generalizar
razonablemente sin regularización explícita es, por su parte, consistente con la
observación de Zhang et al.\ \cite{zhang2017understanding} sobre la generalización
implícita de las redes neuronales profundas.

Sin embargo, hay que señalar algunas limitaciones. La arquitectura
empleada es pequeña y las conclusiones no deben extrapolarse directamente a modelos más
profundos. La métrica \textit{unique\_pctg} depende de la resolución de discretización,
por lo que los valores absolutos no son comparables entre configuraciones distintas,
aunque las tendencias cualitativas sí lo son. Los experimentos estudian cada técnica de
forma aislada; en la práctica, las combinaciones pueden producir efectos no lineales
\cite{li2019understanding}. Por último, el número de repeticiones es reducido, lo que
obliga a interpretar los resultados como tendencias antes que como estimaciones
estadísticamente precisas.
